% BHU PYQ -- Manually transcribed & error-corrected
% Paper: BCHI-327: Labour Welfare and Social Security
% Exam: B.Com. (Hons.) (Semester VI) Examination, 2015-16

\documentclass[11pt,a4paper]{article}
\usepackage[a4paper,top=2.5cm,bottom=2.5cm,left=2.8cm,right=2.8cm,headheight=14pt]{geometry}
\usepackage{amsmath,amssymb}
\usepackage[T1]{fontenc}
\usepackage[utf8]{inputenc}
\usepackage{lmodern,microtype,fancyhdr,enumitem,hyperref,lastpage,parskip}

\hypersetup{
    colorlinks=true,
    urlcolor=black,
    linkcolor=black,
    pdftitle={BCHI-327: Labour Welfare and Social Security -- BHU B.Com. (Hons.) Sem VI 2015-16}
}

\pagestyle{fancy}
\fancyhf{}
\renewcommand{\headrulewidth}{0.4pt}
\renewcommand{\footrulewidth}{0.4pt}
\fancyhead[L]{\small   \textit{Banaras Hindu University}}
\fancyhead[R]{\small   \textit{BCHI-327: Labour Welfare and Social Security}}
\fancyfoot[C]{\small Page  \thepage\ of \pageref{LastPage}}


\newlist{parts}{enumerate}{2}
\setlist[parts,1]{label=(\alph*),leftmargin=*,itemsep=5pt,topsep=3pt}
\setlist[parts,2]{label=(\roman*),leftmargin=*,itemsep=3pt,topsep=2pt}

\newcommand{\pts}[1]{\hfill   \textit{[#1]}}


\begin{document}


\begin{center}
  {\large\bfseries BANARAS HINDU UNIVERSITY}\\[2pt]
  {\normalsize B.Com. (Hons.) --- Semester VI Examination, 2015-16}\\[6pt]
  \rule{0.8\linewidth}{0.5pt}\\[6pt]
  {\large\bfseries Paper No. BCHI-327: Labour Welfare and Social Security}\\[4pt]
  {\normalsize Time: 3 Hours\hfill Maximum Marks: 70}
\end{center}

\rule{\linewidth}{0.4pt}

\smallskip



\noindent   \textit{Instructions: Answer all questions. All questions carry equal marks (14 marks each).}

\bigskip




\noindent   \textbf{Question 1.} \\
``Welfare work intelligently concerned and generally administered in the long run repays the employers.'' Discuss.\pts{14}

\centerline{   \textbf{OR}}

``Employer is bound to implement statutory provisions.'' Comment.\pts{14}

\medskip\rule{\linewidth}{0.2pt}




\noindent   \textbf{Question 2.} \\
Explain the main provisions of the Factories Act, 1948 with regard to labour welfare.\pts{14}

\centerline{   \textbf{OR}}

Discuss the main provisions of safety as given in the Factories Act, 1948.\pts{14}

\medskip\rule{\linewidth}{0.2pt}

\pagebreak




\noindent   \textbf{Question 3.} \\
What do you mean by Social Assistance and Social Insurance? Differentiate between the two.\pts{14}

\centerline{   \textbf{OR}}

What is Social Security? Point out its economic significance in the present era.\pts{14}

\medskip\rule{\linewidth}{0.2pt}




\noindent   \textbf{Question 4.} \\
Is the International Labour Organization a boon or bane for Indian labourers? Give your arguments in favour and against it.\pts{14}

\centerline{   \textbf{OR}}

Discuss the functions of the International Labour Organization.\pts{14}

\medskip\rule{\linewidth}{0.2pt}




\noindent   \textbf{Question 5.} \\
Discuss the salient features of the Maternity Benefit Act, 1961.\pts{14}

\centerline{   \textbf{OR}}

Discuss the main provisions of the Employees' State Insurance Act, 1948.\pts{14}

\vfill
\rule{\linewidth}{0.4pt}

\begin{center}\small
\href{https://bhu-exam.vercel.app/}{Click here to visit our website}\\[4pt]
\href{https://me-aryan.vercel.app/}{Click here to visit our contributor}\\[4pt]
\href{https://quashan-me.vercel.app/}{Click here to visit our contributor}
\end{center}

\end{document}
